\documentclass[../../../include/open-logic-section]{subfiles}

\begin{document}
\olfileid[fr]{sth}{ord-arithmetic}{intro}

\olsection{Introduction}

Dans le \olref[ordinals][]{chap}, nous avons développé une théorie
des nombres ordinaux. Nous avons vu dans le
\olref[spine][]{chap} que les ordinaux peuvent être considérés
comme l'ossature autour de laquelle se construit le reste
de la hiérarchie. Mais ce n'est pas leur seul rôle : il faut
aussi faire de l'arithmétique ordinale.

Nous l'avions déjà évoquée dans la
\olref[ordinals][idea]{sec}, lorsque nous parlions de
$\omega$, $\omega+1$ et $\omega+\omega$. Nous nous
exprimions alors de manière informelle ; le moment est venu
de préciser les choses. Signalons toutefois qu'il y a peu de
philosophie dans ce chapitre : surtout des développements
techniques, accompagnés d'une remarque (modérément)
intéressante sur la possibilité de faire la même chose de
deux manières différentes.

\end{document}
